\scanpage{126}
de traiter le cas \unsure{général} de coefficients non cts sur une courbe algébrique, en donnant notamment explicitement les termes locaux.

Une \inserted{petite} difficulté \inserted{technique} préliminaire dans la formulation, qui se pose dans tous les cas, est que en coh. étale nous disposons à priori de coefficients $A=\mathbb Z/n\mathbb Z$ \inserted{p.ex. $A=\mathbb Z/\ell^\nu\mathbb Z$} qui ne \uline{sont pas des corps} [Il ne serait pas admissible de se limiter à des corps \inserted{tels} $\mathbb Z/\ell\mathbb Z$, car on veut des résultats plus précis que ``le nb des pts fixes mod $\ell$'' --- aussi pour le cas il suffirait alors de prendre la vieille coh. des faisceaux \unsure{cts|ét} !].

Or alors les $H^i(X,F_\bullet)$ ne sont plus des modules libres sur $A$, et les $\operatorname{Tr}f^{(i)}$ ne sont pas définis.

Une méthode, indiquée dans mon exposé Bourbaki, de tourner la difficulté, est de passer de prime abord sur la limite \inserted{projective} des \unsure{cohomologies} à coeff dans des faisceaux comme $\mathbb Z/\ell^\nu\mathbb Z$ $(\nu\longrightarrow\infty)$ de façon à trouver \deleted{des} \deleted{\unsure{corps}} \inserted{des} anneaux de coefficients tels $\Lambda=\mathbb Z_\ell$, \deleted{\unsure{et des}} qui sont au moins \uline{intègres} (et \unsure{à} un corps de fractions $\mathbb Q_\ell$ de car. $0$ !) --- alors
\mspage{125}
